\section{Anexo: respostas às anotações da etapa anterior}
\label{sec:anotacoes}

A versão anterior deste estudo, entregue como atividade do Módulo I, recebeu cinco anotações da
professora. A Tabela \ref{tab:anotacoes} indica onde cada uma foi respondida nos resultados preliminares.

\begin{table}[htb]
\centering
\caption{Anotações da avaliação do Módulo I e onde foram respondidas.}
\label{tab:anotacoes}
\small
\begin{tabular}{p{5.2cm}p{9.6cm}}
\toprule
\textbf{Anotação} & \textbf{Resposta} \\
\midrule
``Onde está essa discussão?'', sobre o confronto entre a taxa de ataque e a soroprevalência &
A discussão dos resultados preliminares compara o SIR com o inquérito do próprio Recife
\cite{braga2023}, que substitui o de Salvador usado antes, e calcula o $R_0$ implícito. \\
\addlinespace
``E pode se reinfectar?'', sobre o intervalo entre duas infecções &
O texto passou a explicar que o intervalo de geração se refere a duas pessoas diferentes na cadeia. A
imunidade permanente é justificada com três estudos, e uma extensão com perda de imunidade mostra o que
mudaria. \\
\addlinespace
``Onde está esse valor?'', sobre a taxa de ataque com seis casas decimais &
A Tabela \ref{tab:validacao} mostra os valores simulado e analítico e o erro relativo. \\
\addlinespace
``Faltaram cenários com gamma maior do que beta!!'' &
São dez cenários nos três regimes (Tabela \ref{tab:cenarios}), com figura própria para $\gamma \geq \beta$,
o total de infectados até a extinção e o mapa de todas as combinações de $\beta$ e $\gamma$. \\
\addlinespace
``Onde está essa análise?'', sobre o controle vetorial de 60\% &
A eficácia deixou de ser um valor fixo sem fonte. Foram variados a eficácia e o dia de início no SIR, e o
controle de criadouros foi simulado no SEIR-SEI. \\
\bottomrule
\end{tabular}
\end{table}

Além das anotações, a revisão das fontes corrigiu a autoria do $R_0$ de 3,4, que é de
\citeonline{haider2020}, e a descrição do estudo de \citeonline{magalhaes2017}, feito numa unidade de saúde
de Paulista, e passou a usar uma única fonte para o período infeccioso \cite{manore2014}. Por isso a taxa
de recuperação mudou de 1/7 para 1/6 por dia, e o pico da linha de base passou do dia 44 para o dia 38.
